\documentclass{article}
\begin{document}
\title{The last useful formation notice depends on when an older receipt stops using the radio budget}
\maketitle

\paragraph{Shared law and scope.}
Two formation plans can reach and dwell at their respective authorized
sensors. A fixed inspection requirement is captured at a scheduled cut.
The moving participant may start a common preparation immediately
afterward, then use the received bit to choose the actual branch.
The radio path is the selected pinned SubT high-signal broker/client
and budget checker; its callback, not attempted delivery, supplies the
bit. The dwell and actual time scaling are modeled, not a lidar,
physical radio or autonomous path to a real station. A preinstalled
source convention ties the one-byte payload to this particular
capture and task; identity and clock qualification are not provided
by a byte on their own.

\paragraph{Compared cases.}
The near station captures at .425 after .42 stationary dwell at the
initial position. Its sharp useful-receipt cutoff is in
$[.604725,.604726]$. The farther station at $x=3.9$ captures
at $\sigma-.15$ after station arrival at .2; its sharp cutoff is
in $[.59228403,.59228405]$. Both are from the exact parent model.
They share a broker clock with evenly spaced send slots of $1/1024$;
set motion time zero one second after broker time zero. The resulting
last legal receipt ticks relative to motion zero are respectively 619
and 606. Each successful callback drives a checked rational branch
selector: the common prefix can defer using the bit to the fixed
certified cut (strictly later than the final legal send slot but
strictly within the parent's bound). The test verifies the resulting
piecewise-quadratic paths at all relevant extrema.

An independent already owed report of two bytes arrives before the
motion horizon, at a pinned tick $o<0$ on that same receiving
interface. Neither dropping nor moving it is permitted; both arms
account for the actual callback at $o$. The truthful cue is a
one-byte message from a distinct sender, unavailable before capture.
At cap 2, its first legal arrival is the later of the first
post-capture send slot and $o+1025$: the old debit is still charged
at $o+1024$ by the source's strict expiry test. Consequently the
latest permitted $o$ is $-406$ at the near station and $-419$ at
the farther station. The broker uses exact binary-representable
timestamps for these cuts. At $o=-411$ only the near station works;
at $-421$ both work and at $-401$ neither does. Adding one byte of
receive capacity admits both reports within the declared window in
all tested cases. This is a paid design change, not a zero-cost
scheduler advantage. An earlier arriving old report can cease to
occupy capacity before the otherwise identical later one.

\paragraph{Comparison and residuals.}
The earliest eligible send rule takes the qualified old receipt's
time as provided to the scheduler. Independent exact integer
accounting and exhaustive slot feasibility agree with it.
The weaker fixed-time and inclusive-boundary controls illustrate
the lost service but are not the strongest baseline. If site and
old receive tick each have equal marginals over near/far and
$-421/-411$, the possible success fraction is $1/2$ to $1$ at
cap 2, depending on the joint law. This is a consequence of the
exact success table, not an empirical risk estimate.

The old obligation is simply an already delivered native report
in this fixture; there is no claim that all prior duties terminate
on reception. A weaker scheduler observing only current byte count
would not know the expiration, but a fully informed direct rule
ties. This is a composition check with a real resource accounting
path, not a full simulator end-to-end field case, an independent
mission, or a proof that teams can choose the station after learning
the old receipt's time.
\end{document}
